% !TEX root = ../Main.tex
\chapter{斜面上物体的平衡问题}

斜面上一个物块静止时，沿斜面方向是重力分量 $mg\sin\theta$ 想把它往下拉、静摩擦力 $f$ 顶住它。物块不下滑，靠的就是摩擦力。把临界情形写出来，会牵出一串看似不同、其实是同一件事的现象。

\begin{center}
\begin{tikzpicture}[>=Stealth, scale=1]
    \draw[thick] (0,0) -- (5,2.5);
    \draw[thick] (0,0) -- (5,0);
    \draw (0.7,0) arc (0:27:0.7);
    \node at (0.95,0.18) {\small $\theta$};
    \coordinate (C) at (2.6,1.3);
    \draw[thick, fill=gray!15] (C) ++(207:0.32) -- ++(27:0.64) -- ++(117:0.5) -- ++(27:-0.64) -- cycle;
    \draw[->, thick] (C) ++(0,0.25) -- ++(117:0.95) node[above left] {\small $N$};
    \draw[->, thick] (C) ++(0,0.05) -- ++(27:1.0) node[right] {\small $f$};
    \draw[->, thick] (C) -- ++(0,-1.0) node[below] {\small $mg$};
\end{tikzpicture}
\end{center}

\section{临界条件}

\state{不下滑的临界条件}{
    沿斜面：重力分量 $mg\sin\theta$，最大静摩擦力 $f_{\max}=\mu N=\mu mg\cos\theta$。物块不下滑要求
    \[
    mg\sin\theta\le \mu mg\cos\theta\quad\Longleftrightarrow\quad \mu\ge\tan\theta.
    \]
}

这一个不等式 $\mu\ge\tan\theta$ 同时引出下面两件事：

\begin{itemize}
    \item[$\Rightarrow$] 自锁现象；
    \item[$\Rightarrow$] 沙堆的倾角。
\end{itemize}

\section{引申：沿斜面加外力}

若 $\mu<\tan\theta$，物块自己会下滑，要让它停住（或匀速上推），就得沿斜面加一个力 $F$。

\begin{center}
\begin{tikzpicture}[>=Stealth, scale=1]
    \draw[thick] (0,0) -- (5,2.5);
    \draw[thick] (0,0) -- (5,0);
    \draw (0.7,0) arc (0:27:0.7);
    \node at (0.95,0.18) {\small $\theta$};
    \coordinate (C) at (2.6,1.3);
    \draw[thick, fill=gray!15] (C) ++(207:0.32) -- ++(27:0.64) -- ++(117:0.5) -- ++(27:-0.64) -- cycle;
    \draw[->, very thick, blue] (C) ++(27:0.35) -- ++(27:1.1) node[right] {\small $F$};
\end{tikzpicture}
\end{center}

\section{冲上斜面}

让物块以初速度 $v_0$ 沿斜面向上冲。

\begin{center}
\begin{tikzpicture}[>=Stealth, scale=1]
    \draw[thick] (0,0) -- (5,2.5);
    \draw[thick] (0,0) -- (5,0);
    \draw (0.7,0) arc (0:27:0.7);
    \node at (0.95,0.18) {\small $\theta$};
    \coordinate (C) at (1.8,0.9);
    \draw[thick, fill=gray!15] (C) ++(207:0.3) -- ++(27:0.6) -- ++(117:0.46) -- ++(27:-0.6) -- cycle;
    \draw[->, very thick, blue] (C) ++(27:0.35) -- ++(27:1.05) node[right] {\small $v_0$};
\end{tikzpicture}
\end{center}

\rmk{注意 $\mu$ 与 $\tan\theta$ 的关系。}
